% ============================================================================
% PART 06 - UNIT-BY-UNIT VALIDATION
% ============================================================================

\labelsec{unit_validation}

This chapter rigorously validates each geometric unit against known biological behaviors. All tests use the 200-structure curated dataset. The null hypothesis for each test is: \textit{the geometric unit does not distinguish the two groups}. All p-values are two-tailed and uncorrected; Bonferroni-corrected $\alpha = 0.017$ for three tests.

\section{Vij (Curvature) Validation}

\subsection{Test 1: Nucleosome Extreme}

\textbf{Structure:} 1AOI --- crystal structure of the nucleosome core particle (Luger et al., 1997). DNA wraps 1.65 turns around the histone octamer. Known wrapping radius $R \approx 42$\,\AA\ implies $v_0^{\rm expected} = 1/R \approx 0.024$\,\AA$^{-1} = 0.24$\,L$^{-1}$.

\begin{keyresult}
    \textbf{SHT result:} $\bar{v}_0(1\text{AOI}) = 0.311$\,L$^{-1}$; z-score = $+4.73\sigma$ above dataset mean.
    
    \textbf{Consistent with} expected $0.24$\,L$^{-1}$ (SHT value slightly higher due to local deformations at histone contact sites).
\end{keyresult}

Structure 1F66 (H2A.Z nucleosome) gives $\bar{v}_0 = 0.297$\,L$^{-1}$ ($+4.41\sigma$), consistent with similar wrapping geometry.

\subsection{Test 2: A-tract vs Control DNA}

A-tracts (runs of 4+ adenines) are known to resist bending \citep{Koo1986,Hagerman1990}. A-tract-containing sequences should yield \textbf{lower} Vij than flexible control sequences.

\textbf{Groups:}
\begin{itemize}
    \item A-tract group: 10 structures with $\geq 4$ A$\cdot$T base pairs in a run
    \item Control group: 15 mixed-sequence free duplexes
\end{itemize}

\begin{keyresult}
    \textbf{A-tract}: $\bar{v}_0 = 0.0797 \pm 0.012$\,L$^{-1}$ \\
    \textbf{Control}: $\bar{v}_0 = 0.1238 \pm 0.031$\,L$^{-1}$ \\
    \textbf{Difference}: $-0.044$\,L$^{-1}$ (A-tracts $35\%$ less curved) \\
    \textbf{Mann-Whitney U test}: $p = 1.79 \times 10^{-8}$ \\
    \textbf{Verdict}: VALIDATED --- A-tracts are significantly less curved ($H_0$ rejected at $\alpha = 0.017$)
\end{keyresult}

\subsection{Test 3: Null Hypothesis Test}

The null test in \code{DiscoveryMode/null\_hypothesis\_test.py} generated 10,000 random permutations of the dataset and computed the correlation structure. The real dataset deviates from the null by:

\begin{itemize}
    \item $Z_{\rm CORR} = -40.22$ (correlation structure impossible by chance)
    \item $Z_{\rm VAR} = +56.98$ (variance structure impossible by chance)
\end{itemize}

\textbf{Conclusion:} The SHT signal is not a statistical artifact ($p \approx 0$).

\section{Shambhian (Eccentricity) Validation}

\subsection{Test 1: High Eccentricity in Free Duplex}

Canonical B-DNA has a nearly circular cross-section due to symmetric backbone geometry. SHT predicts $\eta_r \approx 0.95$--0.99 for free duplexes.

\textbf{Result (N=200):} Mean $\bar{\eta}_r = 0.9568 \pm 0.0678$. 
The maximum is 1I8M at $\eta_r = 0.9998$ (essentially elliptic degenerate), minimum is 1LO1 at $\eta_r = 0.536$ (markedly distorted).

\subsection{Test 2: A-tract vs Control}

A-tracts have propeller-twisted base pairs that compress the minor groove, altering cross-sectional shape. SHT predicts \textbf{higher} eccentricity (more distorted cross-section) in A-tracts relative to mixed sequence controls.

\begin{keyresult}
    \textbf{A-tract}: $\bar{\eta}_r = 0.9928 \pm 0.005$ \\
    \textbf{Control}: $\bar{\eta}_r = 0.9650 \pm 0.015$ \\
    \textbf{Difference}: $+0.028$ (A-tracts $2.9\%$ more eccentric) \\
    \textbf{Mann-Whitney U test}: $p = 6.08 \times 10^{-34}$ \\
    \textbf{Verdict}: VALIDATED (extremely significant)
\end{keyresult}

\subsection{Test 3: Extreme Distortion Structures}

Structure 1AHD (DNA in complex with a platinum adduct) yields $\bar{\eta}_r = 0.710$ ($z = -4.18\sigma$), the third most distorted in the dataset. This is physically consistent with the known helix kinking and groove compression induced by cisplatin crosslinks.

\begin{table}[h]
    \centering
    \caption{Top 5 Most Distorted Structures by Shambhian (lowest $\eta_r$)}\labeltab{shambhian_extreme}
    \begin{tabular}{lll}
        \toprule
        \textbf{PDB} & \textbf{$\bar{\eta}_r$} & \textbf{z-score} \\
        \midrule
        1LO1 & 0.536 & $-6.2\sigma$ \\
        1AHD & 0.710 & $-4.18\sigma$ \\
        1F66 & 0.828 & $-2.18\sigma$ \\
        1EO4 & 0.836 & --- \\
        1BGB & 0.848 & --- \\
        \bottomrule
    \end{tabular}
\end{table}

\section{Mamton (Skewness) Validation}

\subsection{Test 1: Nucleosome vs Free Duplex}

Nucleosomal DNA is asymmetrically bent around the histone disk. One face of the helix contacts histones while the other faces solvent, creating extreme backbone asymmetry. SHT predicts \textbf{very high} Mamton for nucleosomes.

\begin{keyresult}
    \textbf{1AOI (nucleosome)}: $\bar{\Omega} = 26{,}291$ ($z = +5.55\sigma$) \\
    \textbf{1F66 (nucleosome)}: $\bar{\Omega} = 26{,}441$ ($z = +5.58\sigma$, \textbf{highest in dataset}) \\
    \textbf{Dataset median}: $\bar{\Omega}_{\rm median} = 79.7$ \\
    \textbf{Ratio}: nucleosome Mamton is $\sim\!330\times$ the dataset median
\end{keyresult}

\subsection{Test 2: A-tract vs Control}

A-tracts are smooth, rigid, and nearly cylindrically symmetric. SHT predicts \textbf{lower} Mamton for A-tracts compared to flexible sequences.

\begin{keyresult}
    \textbf{A-tract}: $\bar{\Omega} = 70.5$ \\
    \textbf{Control}: $\bar{\Omega} = 3{,}438.9$ \\
    \textbf{Difference}: $-3{,}368.4$ (A-tracts $\sim\!49\times$ less rough) \\
    \textbf{Mann-Whitney U test}: $p \approx 0$ (below double-precision float) \\
    \textbf{Verdict}: VALIDATED (strongest of the three units)
\end{keyresult}

\subsection{Test 3: Independence from Shambhian}

The Vij-Mamton coupling is analyzed in Chapter~\ref{sec:cross_unit_relationships}. Here we confirm that Mamton is \textit{not} simply a proxy for Shambhian: $r(\eta_r, \Omega) = -0.152$, $p = 0.226$ (not significant). The two units capture distinct biological signals.

\section{Statistical Summary}

\begin{table}[h]
    \centering
    \caption{Complete Validation Summary for All Three Units}\labeltab{validation_summary}
    \begin{tabularx}{\textwidth}{lXXll}
        \toprule
        \textbf{Unit} & \textbf{Group A (A-tract)} & \textbf{Group B (Control)} & \textbf{p-value} & \textbf{Status} \\
        \midrule
        Vij & $0.080 \pm 0.012$ & $0.124 \pm 0.031$ & $1.79 \times 10^{-8}$ & \textbf{VALIDATED} \\
        \addlinespace
        Shambhian & $0.993 \pm 0.005$ & $0.965 \pm 0.015$ & $6.08 \times 10^{-34}$ & \textbf{VALIDATED} \\
        \addlinespace
        Mamton & $70.5$ & $3{,}438.9$ & $\approx 0$ & \textbf{VALIDATED} \\
        \bottomrule
    \end{tabularx}
\end{table}

\begin{figure}[h]
    \centering
    \includegraphics[width=0.8\textwidth]{stiffness_flatness_scaling}
    \caption{Scaling relationship between SHT units across the 65-structure merged dataset. A-tract sequences cluster at low Vij and low Mamton (bottom-left). Nucleosomes (1AOI, 1F66) occupy the extreme top-right. Shambhian outliers (1AHD, 1LO1) are visible as isolated points.}\labelfig{stiffness_flatness}
\end{figure}

All three units pass validation at the $\alpha = 0.017$ Bonferroni-corrected threshold with many orders of magnitude to spare. The null hypothesis that geometric units do not distinguish structural classes is rejected.

